Prior work finds that relations expressed more often in pretraining data are more
likely to admit linear decoders in language-model representations. Does exposure
frequency itself determine this geometry? We separate three properties that are
often conflated: recalling a relation, representing it with an affine map, and
executing it with affine computation. We then construct two controlled synthetic
settings. In the first, a one-layer transformer sees each attribute-defined
relation fact 16 times. It reaches \ExpOneSemanticHeldout{} recall on subjects absent from relation
training, while an ordinary least-squares (OLS) map fitted on trained subjects
selects the correct held-out object with \ExpOneSemanticProbe{} accuracy. A frequency-matched random
relation remains at chance under the same probe. In the second, an eight-layer
GPT-style transformer memorizes a random symmetric relation sampled ten times
as often as each of ten attribute-equality controls during joint training.
Random-relation accuracy is
\ExpTwoRandom{}, yet the best overdetermined affine fit on all learned entity rows explains
only \ExpTwoRandomRsq{} of target variance, compared with a mean per-seed minimum of \ExpTwoStructuredRsq{} for
the structured controls. Component interventions show that the forward computation requires
content-dependent attention and GELU MLPs. These results give controlled
counterexamples to a frequency-only account. An aggregation of 5,388 LMEnt-1B
predictions likewise finds mixed associations between corpus frequency and
affine-decoder faithfulness.
